\section{Conclusion}
\label{sec:conclusion}

This survey has catalogued the compression techniques that published methods have used to deploy deep stereo matchers on edge hardware. We grouped them into seven compression families (backbone substitution, cost-volume compression, iterative-loop compression, knowledge distillation, architectural compression, adaptive compute, and neural architecture search), described the mechanism of each, and summarised 41 methods on a common set of axes (KITTI 2015 D1-all, Scene Flow EPE, parameter count, and reported latency) in Table~\ref{tab:efficient_comparison} and Fig.~\ref{fig:pareto_kitti}. Alongside the compression literature we reviewed the cross-domain generalisation techniques that are relevant when a compressed model is deployed outside its training distribution, and the limited body of work that reports latency on an actual embedded accelerator.

A few observations are worth emphasising as a summary. The foundation-model paradigm (DEFOM-Stereo~\cite{jiang2025defom}, FoundationStereo~\cite{wen2025foundationstereo}, MonSter~\cite{cheng2025monster}, Stereo-Anywhere~\cite{bartolomei2025stereoanywhere}) currently sets the upper envelope of stereo accuracy, at parameter counts and inference latencies that are at least an order of magnitude above typical embedded budgets. The seven compression families identified in Section~\ref{sec:taxonomy} appear to be largely orthogonal, and the 2024--2026 wave of efficient-stereo methods (Pip-Stereo~\cite{zheng2026pipstereo}, Distill-then-Prune~\cite{pan2024dtp}, Fast-FoundationStereo~\cite{wen2026fastfoundationstereo}, BANet~\cite{xu2025banet}, LightStereo~\cite{guo2025lightstereo}, GGEV~\cite{liu2026ggev}) commonly combines two of them. Cross-domain generalisation, which Pip-Stereo's DrivingStereo measurements showed is non-trivial for non-iterative methods, appears to be largely a training-recipe issue and therefore composes naturally with any backbone-level compression.

Two methodological points recurred across the surveyed literature and may be useful for future evaluations: reporting latency on a fixed embedded reference platform alongside the desktop-GPU number, and reporting cross-domain accuracy on Middlebury, ETH3D, and DrivingStereo weather subsets in addition to the source-domain KITTI number. Both are inexpensive to add and would substantially improve cross-paper comparability of efficient-stereo work.

\section*{Acknowledgments}
This review draws on the analyses of approximately one hundred deep-stereo publications conducted as part of a broader research project on edge stereo matching. The authors gratefully acknowledge the open availability of source code and pre-trained weights released by the authors of the methods reviewed; without this openness a survey of this depth would not be possible.
